% !TEX root = ./index.tex

\maketitle
\newpage

\section*{Acknowledgement}

I discussed the assessment questions with friends and obtained feedback on my general approach to some parts.
This feedback informed me of mistakes that I may have made, as well as providing alternative perspectives on the method of proof.

These discussions especially helped with understanding Question (2) part (f) and Question (3).

\section*{Answers}

\begin{enumerate}

        % 1.
    \item In the scope of this question, $i$ is taken to be the unit imaginary number.

        \begin{enumerate}

                % 1.a
            \item The binomial theorem \parencite{binomial-theorem} states that for $x, y \in \mathbb{C}, n \in \mathbb{Z}$,
                \[ (x + y)^n = \sum_{k = 0}^n {\binom{n}{k} x^k y^{n - k}}. \]
                where $\binom{n}{k}$ is defined as
                \[ \binom{n}{k} = \frac{n!}{k! (n - k)!}. \]

                \medskip

                Set $x = 1$, $y = ic$, and $n = 5$ to get
                \begin{align*}
                    (1 + ic)^5 &= \binom{5}{0} (ic)^0 + \binom{5}{1} (ic)^1 + \ldots + \binom{5}{5} (ic)^5 \\
                    &= 1 + 5(ic) + 10(ic)^2 + 10(ic)^3 + 5(ic)^4 + (ic)^5 \\
                    &= 1 + 5ci - 10c^2 - 10c^3 i + 5c^4 + c^5 i \\
                    &= (5c^4 - 10c^2 + 1) + c(c^4 - 10c^2 + 5)i.
                \end{align*}

                \bigskip

                % 1.b
            \item
                \begin{proposition} \label{prop:1}
                    $c^4 - 10c^2 + 5 = 0$.
                \end{proposition}

                \begin{proof}
                    It is given that $(1 + ic)^5$ is real, which means that the imaginary component of $(1 + ic)^5$ must be zero, thus \[ c(c^4 - 10c^2 + 5) = 0. \]
                    But c is given to be non-zero, so $c^4 - 10c^2 + 5 = 0$.
                \end{proof}

                \bigskip

                % 1.c
            \item We have
                \begin{flalign*}
                    && c^4 - 10c^2 + 25 &= 20, && \\
                    \implies && (c^2 - 5)^2 &= 20, \\
                    \implies && c^2 - 5 &= \pm \sqrt{20} \\
                    && &= \pm 2 \sqrt{5} \\
                    \text{thus} && c^2 &= 5 + 2\sqrt{5} \quad \text{or} \quad 5 - 2\sqrt{5}.
                \end{flalign*}

                Hence there are four real solutions for $c$:
                \begin{equation}
                    \begin{split} \label{eq:1}
                        c &= \sqrt{5 + 2\sqrt{5}} \quad \text{or} \quad \sqrt{5 - 2\sqrt{5}} \quad \text{or} \\
                        & \quad -\sqrt{5 + 2\sqrt{5}} \quad \text{or} \quad -\sqrt{5 - 2\sqrt{5}}.
                    \end{split}
                \end{equation}

                \bigskip

                % 1.d
            \item Assume that $1 + ic = r e^{i \theta}$, where $r, \theta \in \mathbb{R}$.

                \begin{proposition} \label{prop:2}
                    The smallest positive value of $\theta$ is $\frac{\pi}{5}$.
                \end{proposition}

                \begin{proof}
                    By de Moivre's Theorem,
                    \begin{align*}
                        (1 + ic)^5 &= (r e^{i\theta})^5 \\
                        &= r^5 e^{i(5\theta)}.
                    \end{align*}

                    Because $(1 + ic)^5$ is a real number, the argument $5\theta$ must be of the form $\pi k$ for $k \in \mathbb{Z}$.
                    Thus the positive values of $\theta$ are
                    \[ \theta = \frac{\pi}{5}, \: \frac{2\pi}{5}, \, \ldots, \, \frac{\pi k}{5}. \]

                    It follows that the smallest positive value of $\theta$ is $\frac{\pi}{5}$.
                \end{proof}

                \begin{corollary}
                    The function $\theta \mapsto \tan \theta$ is increasing in the interval $[0, \frac{\pi}{2})$, so the minimum positive value of $\tan \theta$ is $\tan \frac{\pi}{5}$.
                \end{corollary}

                \bigskip

                % 1.e
            \item
                \begin{proposition} \label{prop:3}
                    $\tan{\frac{\pi}{5}} = \sqrt{5 - 2 \sqrt{5}}$.
                \end{proposition}

                \begin{proof}
                    We express $r e^{i\theta}$ in the form $r(\cos \theta + i \sin \theta)$ and write
                    \begin{spreadlines}{\baselineskip}
                        \begin{align*}
                            r e^{i \theta} &= r \cos \theta \, \left(1 + i \, \frac{\sin \theta}{\cos \theta}\right) \\
                            &= r \cos \theta \: (1 + i \tan \theta).
                        \end{align*}
                    \end{spreadlines}

                    But $r e^{i\theta} = 1 + ic$, so by equating real components, $r \cos \theta = 1$. Thus
                    \begin{flalign}
                        && r \cos \theta \: (1 + i \tan \theta) &= 1 + ic, \nonumber && \\
                        \text{so} && 1 + i \tan \theta &= 1 + ic, && \text{($r \cos \theta = 1$)} \nonumber \\
                        \text{hence} && \tan \theta &= c. \label{eq:2}
                    \end{flalign}

                    \medskip

                    Note that $\theta$ and $c$ can take multiple values.
                    There are four unique non-zero values of $\tan \theta$ corresponding to
                    $\theta = \frac{\pi}{5}, \, \frac{2\pi}{5}, \, \frac{3\pi}{5}, \, \frac{4\pi}{5}$,
                    and $c$ can be one of the four possible values in \cref{eq:1}.

                    The numbers of possible values for each variable are equal, so for each value of $\theta$ there is a unique value of $c$ that satisfies \cref{eq:2}, and vice versa.
                    This implies that the minimum positive value of $\tan \theta$ is equal to the minimum positive value of $c$.

                    \medskip

                    But the minimum positive value of $c$ is $\sqrt{5 - 2 \sqrt{5}}$.
                    Therefore, due to the corollary of \cref{prop:2},
                    \[
                        \tan \left(\frac{\pi}{5}\right) = \sqrt{5 - 2\sqrt{5}}. \qedhere
                    \]
                \end{proof}

        \end{enumerate}

        \bigskip

        % 2.
    \item
        \begin{enumerate}

                % 2.a
            \item Observe the domains of each component of $f$:
                \begin{itemize}
                    \item $\frac{x^2 - 1}{x^2 + 1}$ is defined for all $x \in \mathbb{R}$ because the denominator is always positive.

                    \item $\cos x$ is defined for all $x \in \mathbb{R}$.

                    \item $\frac{1}{\sqrt{x}}$ is undefined when $x \le 0$.
                \end{itemize}
                The natural domain $D$ in $\mathbb{R}$ is thus $\{x \in \mathbb{R} \; | \; x > 0\} = (0, \infty)$.

                \medskip

                For the following proposition, recall the definition of injectivity \\ \parencite{injective-function}:
                \begin{definition} \label{def:injective}
                    A function $x \mapsto f(x)$ is \emph{injective} if, for all $x_0, x_1$ in the domain of $f$,
                    \[ f(x_0) = f(x_1) \implies x_0 = x_1. \]
                \end{definition}

                \begin{proposition} \label{prop:4}
                    $f$ is non-injective.
                \end{proposition}

                \begin{proof}
                    Define $f: D \rightarrow \mathbb{R}$ = $g \cdot h$, where
                    \begin{align*}
                        g(x) &= \frac{x^2 - 1}{x^2 + 1}, \\[2\jot]
                        h(x) &= \cos\left(\frac{1}{\sqrt{x}}\right).
                    \end{align*}

                    To prove that $f$ is non-injective, we will show that $g(a)$ and $h(b)$ are equal to zero for some non-equal $a, b \in D$, which implies that $f(a) = f(b) = 0$.
                    The proof then follows from the contrapositive of \cref{def:injective}.

                    \medskip

                    Let $a = 1$ and $b = \frac{4}{\pi^2}$.
                    Note that $a, b \in D$ because $a$ and $b$ are both positive.

                    \smallskip

                    We evaluate $g(a)$ and $h(b)$ to get
                    \[
                        g(a) = \frac{a^2 - 1}{a^2 + 1} = \frac{1^2 - 1}{1^2 + 1} = \frac{0}{2} = 0,
                    \]
                    and
                    \begin{spreadlines}{\baselineskip}
                        \begin{align*}
                            h(b) &= \cos\left(\frac{1}{\sqrt{b}}\right) \\
                            &= \cos\left(\frac{1}{\sqrt{\left(\frac{4}{\pi^2}\right)}}\right) \\
                            &= \cos\left(\pm \frac{\pi}{2}\right) \\
                            &= 0.
                        \end{align*}
                    \end{spreadlines}
                    Thus,
                    \begin{align*}
                        f(a) &= g(a) \cdot h(a) = 0 \cdot h(a) = 0, \\
                        f(b) &= g(b) \cdot h(b) = g(b) \cdot 0 = 0.
                    \end{align*}

                    \smallskip

                    The contrapositive of \cref{def:injective} states that for all $x_0, x_1 \in D$,
                    \[
                        x_0 \ne x_1 \implies f(x_0) \ne f(x_1).
                    \]
                    But we have constructed $a, b \in D$ such that $a \ne b$ and $f(a) = f(b) = 0$. This contradicts \cref{def:injective}, hence $f$ is non-injective.

                    \medskip

                \end{proof}

                \break

                % 2.b
            \item
                Let $E = [1, \infty)$ and redefine $f: E \rightarrow \mathbb{R} = g \cdot h$.

                \begin{proposition} \label{prop:5}
                    $f$ is injective.
                \end{proposition}

                \begin{proof}

                    We will show that $g$ and $h$ are both strictly increasing over the interval $E$;
                    because $f$ is jointly proportional to both $g$ and $h$, it also must be strictly increasing and is thus injective.

                    Recall the definitions of \emph{strictly increasing} and \emph{strictly decreasing} \parencite{monotonic-function}:
                    \begin{definition} \label{def:strictly-increasing}
                        A function $x \mapsto f(x)$ is \emph{strictly increasing} in the set $E$ if for all $x_0, x_1 \in E$,
                        \[ x_0 < x_1 \implies f(x_0) < f(x_1). \]
                    \end{definition}

                    \begin{definition} \label{def:strictly-decreasing}
                        A function $x \mapsto f(x)$ is \emph{strictly decreasing} in the set $E$ if for all $x_0, x_1 \in E$,
                        \[ x_0 < x_1 \implies f(x_0) > f(x_1). \]
                    \end{definition}

                    \smallskip

                    Suppose that $x_0 < x_1$ for some $x_0, x_1 \in E$. Applying \cref{def:strictly-increasing} to $g$, we have
                    \begin{spreadlines}{1.5\baselineskip}
                        \begin{flalign*}
                            && x_0 &< x_1, \\
                            \implies && x_0^2 &< x_1^2, && \text{($x_0, x_1 > 0$)} \\
                            \implies && 2x_0^2 &< 2x_1^2, \\
                            \implies && x_0^2 - x_1^2 &< x_1^2 - x_0^2, \\
                            \implies && x_0^2 x_1^2 + x_0^2 - x_1^2 + 1 &< x_0^2 x_1^2 - x_0^2 + x_1^2 + 1, \\
                            \implies && (x_0^2 - 1)(x_1^2 + 1) &< (x_0^2 + 1)(x_1^2 - 1), \\
                            \implies && \frac{x_0^2 - 1}{x_0^2 + 1} &< \frac{x_1^2 - 1}{x_1^2 + 1}, \\
                            \text{hence} && g(x_0) &< g(x_1).
                        \end{flalign*}
                    \end{spreadlines}

                    The final implication follows because the denominators of the LHS and RHS are always positive.
                    Thus, $g$ is strictly increasing over $E$.

                    \smallskip

                    Now we apply \cref{def:strictly-increasing} to $h$. Again, suppose that $x_0 < x_1$ for some $x_0, x_1 \in E$, so that
                    \begin{spreadlines}{\baselineskip}
                        \begin{flalign*}
                            && x_0 &< x_1, \\
                            \implies && \sqrt{x_0} &< \sqrt{x_1}, && \text{(since $x_0, x_1 \ge 1$)} \\
                            \implies && \frac{1}{\sqrt{x_0}} &> \frac{1}{\sqrt{x_1}}, && \text{($x_0, x_1 \ne 0$)} \\
                            \implies && \cos\left(\frac{1}{\sqrt{x_0}}\right) &< \cos\left(\frac{1}{\sqrt{x_1}}\right).
                        \end{flalign*}
                    \end{spreadlines}

                    The final implication is true because the inner expression $\frac{1}{\sqrt{x}}$ is in $(0, 1]$ for all $x \ge 1$, and the function $x \rightarrow \cos x$ is strictly decreasing in that interval.
                    Since $x_0, x_1 \in E$ implies that $x_0, x_1 \ge 1$, by \cref{def:strictly-decreasing}, the final implication follows.
                    Thus, $h$ is strictly increasing over $E$.

                    \medskip

                    We have shown that $f$ and $g$ are both strictly increasing over the domain $E$, which implies that $f$ is also strictly increasing since it is jointly proportional to both $f$ and $g$.
                    The proof immediately follows from the following lemma:
                    \begin{lemma} \label{lem:1}
                        If the function $f: E \rightarrow \mathbb{R}$ is strictly increasing, then $f$ is injective.
                    \end{lemma}

                    From \cref{def:strictly-increasing}, for all $x_0, x_1 \in E$,
                    \[ x_0 < x_1 \implies f(x_0) < f(x_1).\]

                    Suppose that $x_0 \ne x_1$ for some $x_0, x_1 \in E$.
                    There are two possibilities: either $x_0 < x_1$ or $x_0 > x_1$.
                    Since $x_0$ and $x_1$ are arbitrary, they can be swapped when necessary to ensure that $x_0 < x_1$.

                    Because of this fact, we can assume that $x_0 \ne x_1 \implies x_0 < x_1$ for any choice of $x_0, x_1 \in E$.
                    After combining this with \cref{def:strictly-increasing}, the implication becomes
                    \[ x_0 \ne x_1 \implies x_0 < x_1 \implies f(x_0) < f(x_1) \implies f(x_0) \ne f(x_1), \]
                    or equivalently
                    \[ x_0 \ne x_1 \implies f(x_0) \ne f(x_1). \]

                    The lemma, and thus the proof, follow from the contrapositive of \cref{def:injective}.
                \end{proof}

                % 2.c
            \item
                \begin{proposition} \label{prop:6}
                    $\lim_{x \rightarrow \infty} f(x) = 1$.
                \end{proposition}

                \begin{proof}
                    To compute $\lim_{x \rightarrow \infty} f(x)$, we first compute $\lim_{x \rightarrow \infty} g(x)$ and $\lim_{x \rightarrow \infty} h(x)$.

                    \smallskip

                    Rearranging $g(x)$,
                    \[
                        g(x) = \frac{x^2 - 1}{x^2 + 1} = \frac{1 - \frac{1}{x^2}}{1 + \frac{1}{x^2}}.
                    \]
                    As $x \rightarrow \infty$, $\frac{1}{x^2} \rightarrow 0$, so $g(x) \rightarrow \frac{1}{1} = 1$.

                    \smallskip

                    As $x \rightarrow \infty$, $\frac{1}{\sqrt{x}} \rightarrow 0$, so $h(x) \rightarrow \cos(0) = 1$.

                    \medskip

                    Recall the following theorem of limits \parencite{limit}:
                    \begin{theorem} \label{thm:limit-product}
                        Suppose that $\phi(x)$ and $\psi(x)$ are real-valued functions that converge as $x \rightarrow c$, for some $c \in \mathbb{R}$. Then
                        \[
                            \lim_{x \rightarrow c} \left[\phi(x) \cdot \psi(x)\right] = \left[\lim_{x \rightarrow c} \phi(x)\right] \left[\lim_{x \rightarrow c} \psi(x)\right].
                        \]
                    \end{theorem}

                    \smallskip

                    Using \cref{thm:limit-product},
                    \[ \lim_{x \rightarrow \infty} f(x) = \left[\lim_{x \rightarrow \infty} g(x)\right] \left[\lim_{x \rightarrow \infty} h(x)\right] = 1 \cdot 1 = 1. \qedhere \]
                \end{proof}

                \begin{remark}
                    From \cref{prop:6}, $f$ has a single horizontal asymptote with the equation $y(x) = 1$.
                \end{remark}

                \bigskip

                % 2.d
            \item
                \begin{proposition} \label{prop:7}
                    $\lim_{x \rightarrow 0^+} f(x)$ does not exist.
                \end{proposition}

                \begin{proof}
                    Suppose that $\lim_{x \rightarrow 0^+} f(x)$ exists and is equal to $L$.

                    Note that $\lim_{x \rightarrow 0^+} g(x) = -1$ since
                    \[ g(0) = \frac{0^2 - 1}{0^2 + 1} = -1. \]

                    From \cref{thm:limit-product},
                    \[ L = \left[\lim_{x \rightarrow 0^+} g(x)\right] \left[\lim_{x \rightarrow 0^+} h(x)\right] = - \left[\lim_{x \rightarrow 0^+} h(x)\right]. \]

                    As $x \rightarrow 0^+$, $\frac{1}{\sqrt{x}} \rightarrow \infty$; the limit is defined because $x$ is always positive.
                    This means that with $y = \frac{1}{\sqrt{x}}$, the expression $\lim_{x \rightarrow 0^+} h(x)$ is equivalent to $\lim_{y \rightarrow \infty} \cos y$, which is undefined as $\cos y$ is periodic.

                    \smallskip

                    It follows that $L$ is also undefined, contradicting the supposition that $\lim_{x \rightarrow 0^+} f(x)$ exists. Thus, $\lim_{x \rightarrow 0^+} f(x)$ does not exist.
                \end{proof}

                \bigskip

                %2.e
            \item Suppose that $n$ is defined on the positive integers, and $\{x_n\}_{n \ge 1}$ is a sequence of positive numbers where
                \[ x_n = \frac{1}{(2 \pi n)^2}. \]

                \begin{proposition} \label{prop:8}
                    $\lim_{n \rightarrow \infty} x_n = 0$ and $\lim_{n \rightarrow \infty} f(x_n) = -1$.
                \end{proposition}

                \begin{proof}
                    First we evaluate $\lim_{n \rightarrow \infty} x_n$.
                    The expression $(2 \pi n)^2$ is strictly increasing as $n$ increases, so we can substitute $y = (2 \pi n)^2$ for $n$ to get
                    \begin{equation} \label{eq:3}
                        \lim_{n \rightarrow \infty} x_n = \lim_{y \rightarrow \infty} \, \frac{1}{y} = 0.
                    \end{equation}

                    To evaluate $\lim_{n \rightarrow \infty} f(x_n)$, we will first find $\lim_{n \rightarrow \infty} g(x_n)$ and $\lim_{n \rightarrow \infty} h(x_n)$.

                    \medskip

                    From \cref{eq:3}, the limit of $g(x_n)$ as $x_n \rightarrow \infty$ is
                    \begin{equation} \label{eq:4}
                        \lim_{n \rightarrow \infty} g(x_n) = \lim_{n \rightarrow \infty} \, \frac{x_n^2 - 1}{x_n^2 + 1} = \frac{0 - 1}{0 + 1} = -1.
                    \end{equation}

                    Simplifying $h(x_n)$ returns a constant:
                    \begin{spreadlines}{\baselineskip}
                        \begin{align}
                            h(x_n) &= \cos\left(\frac{1}{\sqrt{x_n}}\right) \nonumber \\
                            &= \cos\left(\frac{1}{\sqrt{\frac{1}{(2 \pi n)^2}}}\right) \nonumber \\
                            &= \cos(2 \pi n) \nonumber \\
                            &= 1 \quad \text{for all $n$}. \label{eq:5}
                        \end{align}
                    \end{spreadlines}

                    Combining \cref{eq:4} and \cref{eq:5} into \cref{thm:limit-product} gives
                    \[
                        \lim_{n \rightarrow \infty} f(x_n) = \left[\lim_{n \rightarrow \infty} g(x_n)\right] \left[\lim_{x \rightarrow \infty} h(x_n)\right] = (-1)(1) = -1.
                    \]

                    Hence $\lim_{n \rightarrow \infty} x_n = 0$ and $\lim_{n \rightarrow \infty} f(x_n) = -1$.
                \end{proof}

                \begin{remark} \label{remark:1}
                    Since $x_n$ is strictly positive, $\{x_n\}$ approaches zero from the right-hand side.
                \end{remark}

                \bigskip

            \item With $D = (0, \infty)$, redefine $f: D \rightarrow (-1, 1)$.

                Recall the definition of surjectivity \parencite{surjective-function}:
                \begin{definition} \label{def:surjective}
                    The function $f: X \rightarrow Y$ is \emph{surjective} if the image of $f$ is equal to its codomain.
                    That is, for all $y \in Y$, there exists some $x \in X$ such that $f(x) = y$.
                \end{definition}

                \medskip

                \begin{proposition} \label{prop:9}
                    $f$ is surjective.
                \end{proposition}

                \begin{proof}
                    We will use \cref{prop:6} and \cref{prop:8} in conjunction with the Intermediate Value Theorem \parencite{ivt}:
                    \begin{theorem} \label{thm:ivt}
                        Suppose that $a, b \in \mathbb{R}$ and $f: [a, b] \rightarrow \mathbb{R}$ is a continuous function.
                        Then the image of $f$ is a closed interval $I$ where
                        \[ \left[\min(f(a), f(b)), \max(f(a), f(b))\right] \subseteq I. \]
                    \end{theorem}

                    Denote the image of $f$ by $I$, and consider the limits of $f(x)$ as $x$ approaches each bound of $D$.

                    \smallskip

                    From \cref{prop:6}, we have $\lim_{x \rightarrow \infty} f(x) = 1$.
                    In other words, as $x$ approaches the upper bound of $D$, $f(x)$ approaches $1$.
                    This means that $I$ contains $1$ (open).

                    With the sequence $\{x_n\}$ from \cref{prop:8}, we have $\lim_{n \rightarrow \infty} x_n = 0$ and $\lim_{x \rightarrow \infty} f(x_n) = -1$.
                    Note that $x_n \in D$ for all $n$, meaning that the statement $x_n \rightarrow 0$ is equivalent to $x \rightarrow 0$.
                    Although $\lim_{x \rightarrow 0^+} f(x)$ is undefined (\cref{prop:5}), we have $f(x_n) \rightarrow -1$ as $x_n \rightarrow 0^+$ (the lower bound of $D$), hence $I$ contains $-1$ (open).

                    \smallskip

                    So $I$ contains the values $-1$ and $1$ corresponding to the lower and upper bounds of $D$, respectively.
                    Since $f$ is continuous, we can use \cref{thm:ivt} to say that $I$ is a closed interval and that $(-1, 1) \subseteq I$.

                    But the codomain of $f$ is $(-1, 1)$, meaning that every element in the codomain is mapped to by some value in $D$. Hence $f$ is surjective by \cref{def:surjective}.
                \end{proof}

                \bigskip

        \end{enumerate}

        % 3.
    \item Let $1 \le i < j \le 6$ where $i$ and $j$ are integers.
        We want to find the sum of $\vb{u_i} \times \vb{u_j}$ for all combinations of $i, j$.

        To simplify the sum, we make a series of remarks concerning the cross product of $\vb{u_i}$ and $\vb{u_j}$.

        \smallskip

        \begin{remark}
            When summing $\vb{u_i} \times \vb{u_j}$, it is important to consider the direction of each resultant vector.
            For example, by using the right-hand rule, we see that $\vb{u_1} \times \vb{u_2}$ points in the positive $z$ directions, but $\vb{u_1} \times \vb{u_6}$ points in the negative $z$ direction.
        \end{remark}

        \smallskip

        Let $\vb{e_3} = [0, 0, 1]$ denote the unit vector in the positive $z$ direction.
        As a general rule of the previous remark:
        \begin{itemize}
            \item If $j - i = 1$ or  $2$, \\
                \hspace*{1em}then $\vb{u_i} \times \vb{u_j}$ points towards $\vb{e_3}$.
            \item If $j - i = 3$, \\
                \hspace*{1em}then $\vb{u_i} \times \vb{u_j}$ is the zero vector.
            \item If $j - i = 4$ or  $5$, \\
                \hspace*{1em}then $\vb{u_i} \times \vb{u_j}$ points towards $- \vb{e_3}$.
        \end{itemize}

        \smallskip

        \begin{remark}
            For evaluating the cross product, notice that the angle $\theta$ between $\vb{u_i}$ and $\vb{u_j}$ is always a multiple of $\frac{\pi}{3}$: the central angle of the hexagon.
            This multiple happens to be either $(j - i)$ or $6 - (j - i)$.
            To ensure that $\theta$ is non-reflex, set $\theta = \min((j - i), 6 - (j - i)) \, \frac{\pi}{3}$.
        \end{remark}

        Additionally, it is given that $\vb{u_i}, \vb{u_j}$ are unit vectors, so
        \begin{equation} \label{eq:6}
            |\vb{u_i} \times \vb{u_j}| = |\vb{u_i}| |\vb{u_j}| \sin \theta = \sin \theta.
        \end{equation}

        \smallskip

        \begin{remark}
            Because the vectors $\vb{u_1} \ldots \vb{u_6}$ all lie in the $xy$-plane, $\vb{u_i} \times \vb{u_j}$ is always parallel to $\vb{e_3}$, hence the sum $\sum \vb{u_i} \times \vb{u_j}$ for all $1 \le i < j \le 6$ is equal to
            \[ \sum \pm |\vb{u_i} \times \vb{u_j}| \, \vb{e_3}, \]
            where the sign depends on the direction of $\vb{u_i} \times \vb{u_j}$.
        \end{remark}

        \medskip

        In light of the remarks described above, we arrive at the expression
        \begin{spreadlines}{\baselineskip}
            \begin{align*}
                \sum \vb{u_i} \times \vb{u_j} &= (|\vb{u_1} \times \vb{u_2}| + |\vb{u_2} \times \vb{u_3}| + \ldots + |\vb{u_5} \times \vb{u_6}|) \, \vb{e_3} \\
                & \qquad + (|\vb{u_1} \times \vb{u_3}| + |\vb{u_2} \times \vb{u_4}| + |\vb{u_3} \times \vb{u_5}| + |\vb{u_4} \times \vb{u_6}|) \, \vb{e_3} \\
                & \qquad - (|\vb{u_1} \times \vb{u_5}| + |\vb{u_2} \times \vb{u_6}|) \, \vb{e_3} \\
                & \qquad - |\vb{u_1} \times \vb{u_6}| \, \vb{e_3}. \\
                \intertext{By using \cref{eq:6},}
                \sum \vb{u_i} \times \vb{u_j} &= \left(5 \sin\left(\frac{\pi}{3}\right) + 4 \sin\left(\frac{2 \pi}{3}\right) - 2 \sin\left(\frac{2 \pi}{3}\right) - \sin\left(\frac{\pi}{3}\right)\right) \, \vb{e_3} \\
                &= \left(5 \cdot \frac{\sqrt{3}}{2} + 4 \cdot \frac{\sqrt{3}}{2} - 2 \cdot \frac{\sqrt{3}}{2} - \frac{\sqrt{3}}{2}\right) \, \vb{e_3} \\
                &= 3 \sqrt{3} \, \vb{e_3}.
            \end{align*}
        \end{spreadlines}

        Therefore, we have
        \[
            \sum_{\substack{i, j = 1 \\ i < j}}^6 \vb{u_i} \times \vb{u_j} = 3 \sqrt{3} \, \vb{e_3} = [0, 0, 3 \sqrt{3}].
        \]

        \bigskip

        % 4.
    \item
        \begin{enumerate}

                % 4.a
            \item Consider the set of functions
                \[
                    \mathscr{G} = \left\{f: \mathbb{N} \rightarrow \mathbb{R} \: \big| \: f(m) + f(m + 1) = f(m + 2) \; \text{for all} \; m \in \mathbb{N}\right\},
                \]
                together with the pointwise addition and scalar multiplication operations in $\mathbb{R}$: for $f, g \in \mathscr{G}$ and $c \in \mathbb{R}$,
                \begin{align}
                    && (f + g)(m) &:= f(m) + g(m) && \label{eq:7}  \\
                    \text{and} && (cf)(m) &:= cf(m). \label{eq:8}
                \end{align}

                \begin{proposition} \label{prop:10}
                    For all $f, g \in \mathscr{G}$,
                    \[
                        f = g \Longleftrightarrow f(0) = g(0) \; \text{and} \; f(1) = g(1).
                    \]
                \end{proposition}

                \begin{proof}
                    The right implication follows immediately from the fact that the statement $f = g$ means that $f(n) = g(n)$ for all $n \in \mathbb{N}$.

                    \smallskip

                    \begin{note}
                        Since $\mathscr{G}$ is a vector space, there exists a null function $O$ such that $O + \phi = \phi$ for all $\phi \in \mathscr{G}$. Also, for any function $\phi \in \mathscr{G}$, there exists a function $- \phi \in \mathscr{G}$ such that $\phi + (- \phi) = O$.
                    \end{note}

                    To prove the left implication, we will use the Principle of Strong Mathematical Induction to show that
                    \[
                        f(0) = g(0) \; \text{and} \; f(1) = g(1) \implies (f - g)(n) = O(n)
                    \]
                    for all $n \in \mathbb{N}$.

                    \medskip

                    Assume that $f(0) = g(0)$ and $f(1) = g(1)$.

                    Define the expression $f - g$ to be $f + (-g)$. Using the properties of $\mathscr{G}$ described above, we have
                    \begin{align*}
                        f(0) = g(0) &\implies f(0) - g(0) = O(0), \\
                        &\implies (f - g)(0) = O(0) && \text{(\cref{eq:7})} \\[2\jot]
                        f(1) = g(1) &\implies f(1) - g(1) = O(1) \\
                        &\implies (f - g)(1) = O(1). && \text{(\cref{eq:7})}
                    \end{align*}
                    These are the base cases for the induction.

                    Suppose that for some integer $k \ge 0$, the statement $(f - g)(n) = O(n)$ is true for all $n \in \{0, 1, \ldots, k\}$.
                    This means that the statement is true for $n = k - 1$ and $n = k$, that is,
                    \begin{align*}
                        (f - g)(k - 1) &= O(k - 1), \\
                        (f - g)(k) &= O(k).
                    \end{align*}
                    From the definition of $\mathscr{G}$, we make the inductive step:
                    \begin{flalign}
                        && (f - g)(k - 1) + (f - g)(k) &= O(k - 1) + O(k) && \nonumber \\
                        \implies && (f - g)(k + 1) &= O(k + 1).
                    \end{flalign}

                    By Strong Mathematical Induction, it follows from the base cases and inductive step that $(f - g)(n) = O(n)$ for all $n \in \mathbb{N}$.

                    To complete the proof, simply rearrange to get $f(n) = g(n)$.
                    Hence
                    \[
                        f(0) = g(0) \; \text{and} \; f(1) = g(1) \implies f = g. \qedhere
                    \]
                \end{proof}

                \bigskip

                % 4.b
            \item
                \begin{remark}
                    For any function $\Phi \in \mathscr{G}$, $\Phi(0)$ and $\Phi(1)$ behave like the the $x$ and $y$ coordinates of $\Phi$.
                    That is, there exists a function $A: \mathbb{R}^2 \rightarrow \mathscr{G}$ such that $A(\left[\Phi(0), \Phi(1)\right]) = (m \mapsto \Phi(m))$.
                \end{remark}

                To justify this remark, consider successive values of $\Phi(m)$ for some $\Phi \in \mathscr{G}$, starting from $m = 2$:
                \begin{align*}
                    \Phi(2) &= \Phi(0) + \Phi(1), \\[3\jot]
                    \Phi(3) &= \Phi(1) + \Phi(2) \\
                    &= \Phi(0) + (\Phi(0) + \Phi(1)) \\
                    &= \Phi(0) + 2 \cdot \Phi(1), \\[3\jot]
                    \Phi(4) &= \Phi(2) + \Phi(3) \\
                    &= (\Phi(0) + \Phi(1)) + (\Phi(0) + 2 \cdot \Phi(1)) \\
                    &= 2 \cdot \Phi(0) + 3 \cdot \Phi(1), \\[3\jot]
                    \Phi(5) &= \Phi(3) + \Phi(4) \\
                    &= (\Phi(0) + 2 \cdot \Phi(1)) + (2 \cdot \Phi(0) + 3 \cdot \Phi(1)) \\
                    &= 3 \cdot \Phi(0) + 5 \cdot \Phi(1).
                \end{align*}
                Observe that the value of $\Phi(m)$ for any $m \in \mathbb{N}$ can be written as a linear combination of $\Phi(0)$ and $\Phi(1)$.
                This is analogous to how any vector in $\mathbb{R}^2$ can be written as a linear combination of $\vb{e_1} = [1, 0]$ and $\vb{e_2} = [0, 1]$.

                Also, the coefficients of $\Phi(0)$ and $\Phi(1)$ obey a determined pattern, which happens to be the fibonacci sequence defined by the recursive relation $F_m + F_{m + 1} = F_{m + 2}$, with initial values $F_0 = 0, F_1 = 1$.

                For $\Phi(m)$, these coefficients happen to be $F_{m - 1}$ and $F_m$, respectively, so that
                \[ \Phi(m) = F_{m - 1} \, \Phi(0) + F_m \, \Phi(1). \]

                To see that this definition of the function obeys the condition $\Phi(m) + \Phi(m + 1) = \Phi(m + 2)$, we will use mathematical induction.

                \medskip

                \begin{proposition} \label{prop:11}
                    For all functions $\Phi$ in $\mathscr{G}$ and integers $m \ge 1$,
                    \[ \Phi(m) = F_{m - 1} \, \Phi(0) + F_m \, \Phi(1). \]
                \end{proposition}

                \begin{proof}
                    The base cases for the proposition are $m = 1$ and $m = 2$.

                    When $m = 1$,
                    \begin{align*}
                        \text{LHS} &= \Phi(1), \\[2\jot]
                        \text{RHS} &= F_0 \, \Phi(0) + F_1 \, \Phi(1) \\
                        &= 0 \cdot \Phi(0) + 1 \cdot \Phi(1) \\
                        &= \Phi(1) \\
                        &= \text{LHS}.
                    \end{align*}
                    When $m = 2$,
                    \begin{align*}
                        \text{LHS} &= \Phi(2) \\
                        &= \Phi(0) + \Phi(1), \\[2\jot]
                        \text{RHS} &= F_1 \, \Phi(0) + F_2 \, \Phi(1) \\
                        &= F_1 \, \Phi(0) + (F_0 + F_1) \, \Phi(1) \\
                        &= 1 \cdot \Phi(0) + (0 + 1) \, \Phi(1) \\
                        &= \Phi(0) + \Phi(1) \\
                        &= \text{LHS}.
                    \end{align*}

                    Hence the proposition is true for the base cases $m = 1, 2$.

                    Assume that the proposition holds for all $i \in \{1 \ldots k\}$.
                    Making the inductive step $m = k + 1$,
                    \begin{align*}
                        \text{LHS} &= \Phi(k + 1) \\
                        &= \Phi(k - 1) + \Phi(k), \\[2\jot]
                        \text{RHS} &= F_k \, \Phi(0) + F_{k + 1} \, \Phi(1) \\
                        &= (F_{k - 2} + F_{k - 1}) \, \Phi(0) + (F_{k - 1} + F_k) \, \Phi(1) \\
                        &= (F_{k - 2} \, \Phi(0) + F_{k - 1} \, \Phi(1)) + (F_{k - 1} \, \Phi(0) + F_k \, \Phi(1)) \\
                        &= \Phi(k - 1) + \Phi(k) \\
                        &= \text{LHS}.
                    \end{align*}
                    Hence if the proposition is true when $m = i$ for all $i \in \{1 \ldots k\}$, then it is also true when $m = k + 1$.

                    We have proven that the base cases and the inductive step are true; thus, by mathematical induction, the proposition is true for all $m \ge 1$.
                \end{proof}

                \bigskip

                We have that any function $\Phi$ in $\mathscr{G}$ can be described as a linear combination of $\Phi(0)$ and $\Phi(1)$,
                so by using the result from \cref{prop:11}, we obtain a function $A: \mathbb{R}^2 \rightarrow \mathscr{G}$:
                \begin{equation} \label{eq:9}
                    A(\vb{u}) = [m \mapsto \vb{F}(m) \cdot \vb{u}], \quad \vb{u} \in \mathbb{R}^2,
                \end{equation}
                where $\vb{F}(m) = [F_{m - 1}, F_m]$ for $m \ge 1$, the symbol $\cdot$ denotes the dot product of vectors, and $\vb{u} = [\Phi(0), \Phi(1)]$.

                \smallskip

                Note that when substituting $\vb{u} = [\Phi(0), \Phi(1)]$ into the definition of $A$,
                \[
                    A(\vb{u}) = [m \mapsto F_{m - 1} \, \Phi(0) + F_m \, \Phi(1)] = [m \mapsto \Phi(m)].
                \]
                Since $[\Phi(0), \Phi(1)]$ is arbitrary and $\Phi \in \mathscr{G}$, every element in $\mathbb{R}^2$ maps to $\mathscr{G}$ under $T$ --- the domain and codomain of $A$ are $\mathbb{R}^2$ and $\mathscr{G}$, respectively.

                \medskip

                For the following proposition, recall the definition of an isomorphism \parencite{isomorphisms}:
                \begin{definition} \label{def:isomorphic}
                    The vector spaces $V$ and $W$ are \emph{isomorphic} if there exists a function $T: V \rightarrow W$ (called an \emph{isomorphism}) that satisfies the following conditions:
                    \begin{itemize}
                        \item $T(\vb{u} + \vb{v}) = T(\vb{u}) + T(\vb{v})$ for all $\vb{u}, \vb{v} \in V$,
                        \item $T(c \, \vb{u}) = c \, T(\vb{u})$ for all $\vb{u} \in V$ and $c \in \mathbb{R}$,
                        \item $T$ is injective (\cref{def:injective}),
                        \item $T$ is surjective (\cref{def:surjective}).
                    \end{itemize}
                \end{definition}

                \medskip

                \begin{proposition} \label{prop:12}
                    $\mathbb{R}^2$ and $\mathscr{G}$ are isomorphic.
                \end{proposition}

                \begin{proof}
                    Consider the function $A$ from \cref{eq:9}.
                    We divide the main proposition into four claims corresponding to the conditions given in \cref{def:isomorphic}:

                    \begin{claim} \label{subprop:12.1}
                        $A(\vb{u} + \vb{v}) = A(\vb{u}) + A(\vb{v})$ for all $\vb{u}, \vb{v} \in \mathbb{R}^2$.
                    \end{claim}

                    \begin{claim} \label{subprop:12.2}
                        $A(c \, \vb{u}) = c \, A(\vb{u})$ for all $\vb{u} \in \mathbb{R}^2$ and $c \in \mathbb{R}$.
                    \end{claim}

                    \begin{claim} \label{subprop:12.3}
                        $A$ is injective.
                    \end{claim}

                    \begin{claim} \label{subprop:12.4}
                        $A$ is surjective.
                    \end{claim}

                    If each of these claims are true, then $A$ is an isomorphism and the proof immediately follows from \cref{def:isomorphic}.

                    \begin{subproof}
                        For any $\vb{u}, \vb{v} \in \mathbb{R}^2$,
                        \begin{align*}
                            && \text{LHS} &= A(\vb{u} + \vb{v}) \\
                            && &= [m \mapsto \vb{F}(m) \cdot (\vb{u} + \vb{v})] \\
                            && &= [m \mapsto (\vb{F}(m) \cdot \vb{u}) + (\vb{F}(m) \cdot \vb{v})] \\
                            && &= [m \mapsto \vb{F}(m) \cdot \vb{u}] + [n \mapsto \vb{F}(n) \cdot \vb{v}] && \text{(\cref{eq:7})} \\
                            && &= A(\vb{u}) + A(\vb{v}) \\
                            && &= \text{RHS}.
                        \end{align*}
                    \end{subproof}

                    \bigskip

                    \begin{subproof}
                        For any $\vb{u} \in \mathbb{R}^2$, $c \in \mathbb{R}$,
                        \begin{align*}
                            && \text{LHS} &= A(c \, \vb{u}) \\
                            && &= [m \mapsto \vb{F}(m) \cdot (c \, \vb{u})] \\
                            && &= [m \mapsto c \, (\vb{F}(m) \cdot \vb{u})] \\
                            && &= c \, [m \mapsto \vb{F}(m) \cdot \vb{u}] && \text{(\cref{eq:8})} \\
                            && &= c \, A(\vb{u}) \\
                            && &= \text{RHS}.
                        \end{align*}
                    \end{subproof}

                    \bigskip

                    \begin{subproof}
                        From \cref{def:injective}, we want to prove that for all $\vb{u}, \vb{v} \in \mathbb{R}^2$,
                        \[ A(\vb{u}) = A(\vb{v}) \implies \vb{u} = \vb{v}. \]

                        \smallskip

                        For arbitrary $\vb{u}, \vb{v} \in \mathbb{R}^2$, suppose that $A(\vb{u}) = U$ and $A(\vb{v}) = V$, where $U$ and $V$ are in $\mathscr{G}$.

                        The definition of $A$ defines $\vb{u} = [U(0), U(1)]$ and $\vb{v} = [V(0), V(1)]$.
                        But \cref{prop:10} states that if $U = V$, then $U(0) = V(0)$ and $U(1) = V(1)$.
                        Hence if $A(\vb{u}) = A(\vb{v})$, then $U = V$, which implies $\vb{u} = \vb{v}$.
                    \end{subproof}

                    \bigskip

                    \begin{subproof}
                        From \cref{def:surjective}, we want to prove that for every vector $\vb{u} \in \mathbb{R}^2$, there exists some function $f \in \mathscr{G}$ such that $A(\vb{u}) = f$.

                        \smallskip

                        Consider an arbitrary vector $\vb{u} \in \mathbb{R}^2$ and the function $f \in \mathscr{G}$ defined by $[f(0), f(1)] = \vb{u}$.
                        We say that $[f(0), f(1)]$ \emph{defines} the function because of the result from \cref{prop:11},
                        which states that for any $m \in \mathbb{N}$, $f(m)$ can be expressed as a linear combination of $f(0)$ and $f(1)$.

                        \smallskip

                        But $A(\vb{u})$ evaluates to
                        \[
                            (m \mapsto \vb{F}(m) \cdot \vb{u}) = (m \mapsto F_{m - 1} \, f(0) + F_m \, f(1)),
                        \]
                        which, by \cref{prop:11}, is equal to $f$. We have thus constructed a function $f \in \mathscr{G}$ from $\vb{u}$ such that $A(\vb{u}) = f$. Since $\vb{u}$ is arbitrary, we can construct such a function for any element of $\mathbb{R}^2$, which implies that for every $\vb{u} \in \mathbb{R}^2$, there exists a function $f$ in $\mathscr{G}$ such that $A(\vb{u}) = f$. Hence $A$ is surjective by \cref{def:surjective}.
                    \end{subproof}

                    \bigskip

                    It follows from \crefrange{subprop:12.1}{subprop:12.4} that $A$ is an isomorphism of $\mathbb{R}^2$ and $\mathscr{G}$, thereby proving from \cref{def:isomorphic} that $\mathbb{R}^2$ and $\mathscr{G}$ are isomorphic.
                \end{proof}

        \end{enumerate}

\end{enumerate}
